\documentclass[../../../include/open-logic-section]{subfiles}

\begin{document}
\olfileid{sth}{card-arithmetic}{expotough}

\olsection[काही सुलभीकरणे]{संचांक घातांकनाची काही सुलभीकरणे}

$\canonord$ ची व्याख्या थोडी गुंतागुंतीची होती; पण त्यातून संचांक बेरीज आणि गुणाकाराविषयी \olref[simp]{cardplustimesmax} हा उपयोगी निष्कर्ष मिळाला. अनंतपार घातांकन मात्र इतक्या सरळ रीतीने सुलभ करता येत नाही. असे का, हे समजावण्यासाठी सांत प्रकरणातील परिचित नमुन्याचा विस्तार करणाऱ्या निष्कर्षापासून सुरुवात करू (त्याचा पुरावा मात्र बऱ्याच अमूर्त पातळीवर आहे):

\begin{prop}\ollabel{simplecardexpo}
कोणत्याही संचांक $\cardfont{a}, \cardfont{b}, \cardfont{c}$ साठी
$\cardexpo{\cardfont{a}}{\cardfont{b} \cardplus \cardfont{c}} =
\cardexpo{\cardfont{a}}{\cardfont{b}} \cardtimes
\cardexpo{\cardfont{a}}{\cardfont{c}}$ आणि
$\cardexpo{(\cardexpo{\cardfont{a}}{\cardfont{b}})}{\cardfont{c}} =
\cardexpo{\cardfont{a}}{\cardfont{b} \cardtimes \cardfont{c}}$.
\end{prop}

\begin{proof}
पहिल्या विधानासाठी $f \colon
(\cardfont{b}\disjointsum\cardfont{c}) \to \cardfont{a}$ हे फलन पाहा. प्रत्येक
$\beta \in \cardfont{b}$ साठी $f_\cardfont{b}(\beta) = f(\beta, 0)$ आणि प्रत्येक
$\gamma \in \cardfont{c}$ साठी $f_\cardfont{c}(\gamma) = f(\gamma, 1)$ ठरवून त्याचे दोन भाग करा. मग $f \mapsto
\tuple{f_{\cardfont{b}}, f_{\cardfont{c}}}$ हे
$\funfromto{\cardfont{b} \disjointsum \cardfont{c}}{\cardfont{a}} \to
(\funfromto{\cardfont{b}}{\cardfont{a}} \times
\funfromto{\cardfont{c}}{\cardfont{a}})$ यांतील !!a{bijection} आहे.

दुसऱ्या विधानासाठी $f \colon \cardfont{c} \to
(\funfromto{\cardfont{b}}{\cardfont{a}})$ हे फलन पाहा; म्हणजे प्रत्येक $\gamma \in
\cardfont{c}$ साठी $f(\gamma) \colon \cardfont{b} \to
\cardfont{a}$ हे फलन मिळते. आता प्रत्येक $\tuple{\beta, \gamma} \in \cardfont{b} \times \cardfont{c}$ साठी $f^*(\beta, \gamma) = (f(\gamma))(\beta)$ ठरवा.
मग $f \mapsto f^*$ हे
$\funfromto{\cardfont{c}}{(\funfromto{\cardfont{b}}{\cardfont{a}})}
\to \funfromto{\cardfont{b} \times \cardfont{c}}{\cardfont{a}}$ यांतील !!a{bijection} आहे. येथे प्रांत प्रत्यक्ष कार्तीय गुणाकार आहे; त्या प्रांताच्या संचांकानेच संचांक गुणाकार परिभाषित होतो.
\end{proof}

अनंत संचांकांच्या बाबतीत $\cardexpo{\cardfont{a}}{\cardfont{b}}$ सहज मोजता यावा, असे आपल्याला \emph{हवे} आहे. या दिशेने पुढील पाऊल उपयोगी आहे:

\begin{prop}\ollabel{cardexpo2reduct}
$2 \leq \cardfont{a} \leq \cardfont{b}$ आणि $\cardfont{b}$ अनंत असेल, तर $\cardexpo{\cardfont{a}}{\cardfont{b}} =
\cardexpo{2}{\cardfont{b}}$.
\end{prop}

\begin{proof}
\begin{align*}
  \cardexpo{2}{\cardfont{b}} &\leq
  \cardexpo{\cardfont{a}}{\cardfont{b}}\text{, कारण $2 \leq \cardfont{a}$}\\
  &\leq \cardexpo{(\cardexpo{2}{\cardfont{a}})}{\cardfont{b}}
  \text{, \olref[opps]{lem:SizePowerset2Exp} नुसार}\\
  &= \cardexpo{2}{\cardfont{a} \cardtimes \cardfont{b}}
  \text{, \olref{simplecardexpo} नुसार} \\
  &= \cardexpo{2}{\cardfont{b}}
  \text{, \olref[simp]{cardplustimesmax} नुसार}
\end{align*}
\end{proof}

यापुढे आणखी सुलभीकरण मिळेल अशी फारशी अपेक्षा ठेवू नये: \olref[card-arithmetic][opps]{lem:SizePowerset2Exp} नुसार $\cardfont{b} < \cardexpo{2}{\cardfont{b}}$. पण $\cardfont{b} <
\cardfont{a}$ असेल, तर $\cardexpo{\cardfont{a}}{\cardfont{b}}$ बद्दल काय सांगावे, हे यावरून ठरत नाही. अर्थात $\cardfont{b}$ \emph{सांत} असेल, तर काय करायचे ते माहीत आहे.

\begin{prop}
$\cardfont{a}$ अनंत असेल आणि $n \in \omega$ शून्येतर असेल, तर
$\cardexpo{\cardfont{a}}{n} = \cardfont{a}$.
\end{prop}

\begin{proof}
$\cardexpo{\cardfont{a}}{n} = \cardfont{a} \cardtimes \cardfont{a}
\cardtimes \ldots \cardtimes \cardfont{a} = \cardfont{a}$, हे \olref[simp]{cardplustimesmax} नुसार.
\end{proof}
\noindent
याशिवाय, आणखी काही प्रकरणांत $\cardexpo{\cardfont{a}}{\cardfont{b}}$ चे आकारमान ठरवता येते:

\begin{prop}
$2 \leq \cardfont{b} < \cardfont{a} \leq
\cardexpo{2}{\cardfont{b}}$ आणि $\cardfont{b}$ अनंत असेल, तर
$\cardexpo{\cardfont{a}}{\cardfont{b}} = \cardexpo{2}{\cardfont{b}}$.
\end{prop}

\begin{proof}
$\cardexpo{2}{\cardfont{b}}\leq \cardexpo{\cardfont{a}}{\cardfont{b}}
\leq \cardexpo{(\cardexpo{2}{\cardfont{b}})}{\cardfont{b}} =
\cardexpo{2}{\cardfont{b}\cardtimes\cardfont{b}} =
\cardexpo{2}{\cardfont{b}}$; युक्तिवाद \olref{cardexpo2reduct} प्रमाणे.
\end{proof}
\noindent
पण यापुढे गोष्टी अधिक सूक्ष्म होतात.

\end{document}
